% STATUS: written, 20 October 2026.
%
% Most of this is transcribed from docs/mdp-specification.md and
% docs/adr/ADR-0002-flat-declarative-action-space.md. If either of those
% changes, change this with it.

\section{System Model and Methodology}

\subsection{What is being decided}

A cluster has $R$ regions, each holding one database node. The keyspace is
divided into $C$ atomic cells, and a \emph{shard} is a contiguous run of cells.
Each shard has one \emph{primary} copy, which serves writes, and some number of
\emph{replicas}, which hold the same data and can serve reads. The number of
replicas per shard, called the replication factor, is fixed at three throughout
this work and is not something the agent controls.

Three decisions are therefore available at any moment. Which region holds a
given shard's primary. Which regions hold its replicas. Where the boundary
between two shards falls. The first two are allocation and the third is
fragmentation, and this work treats them as one problem because they interact:
splitting a shard is only worth doing if the two halves are then placed
differently.

\subsection{The cost of a placement}

Latency is computed analytically from a fixed matrix of round-trip times
between regions. A read costs the round trip to its nearest copy. A write costs
the round trip to the primary plus the time for a majority of replicas to
acknowledge it, since a write is not durable until enough copies agree. A
transaction touching several shards costs its member operations plus, if those
shards' primaries sit in more than one region, an additional charge for the
two-phase commit protocol that keeps it atomic.

That last term is the mechanism this work is built on. It is what makes placing
together two shards that are used together worth something, and therefore what
makes the fragmentation decision interact with the allocation decision.

\subsection{The workload}

The workload generator produces a stream of requests whose geographic origin
moves. Each of the $C$ cells sits at an angular position on a circle, as do the
regions, and traffic for a cell originates predominantly from whichever regions
sit near its angle. The whole angular field rotates over a daily period, which
is what makes a placement chosen at one time wrong at a later one. Cell
popularity follows a Zipf distribution, so a minority of cells carry most of
the traffic, and localized spikes arrive as a Poisson process, multiplying one
region's traffic for a short burst. The read and write mix and the skew
parameter follow the YCSB conventions~\cite{cooper2010ycsb}, so the workload is
recognizable rather than invented, with geographic origin added as the dimension
YCSB does not model.

Adjacent cells are given correlated but not identical angular positions by a
random walk. This is the detail that makes splitting worthwhile: a wide shard
spans divergent angles and therefore cannot be placed well anywhere, whilst its
two halves can each be placed well.

A single parameter, \emph{volatility}, scales the rotation speed, the spike
rate and the spike size together, so one number moves the environment from
placid to violent. It is the main axis the evaluation sweeps.

\subsection{The decision problem}

The problem is stated as a Markov decision process. The agent acts once per
decision window.

\paragraph{Observation} The state is encoded as a fixed-length vector
describing, for every shard, its size and the regions holding its primary and
replicas; for every region, its share of requests and its storage occupancy;
the recent history of data movement; the matrix of inter-region round-trip
times; and a small set of scalars summarizing current cost.

\paragraph{Actions} The action space is flat and declarative, with 261 actions
at $R=6$ regions and a cap of 20 shards. Four kinds:

\begin{itemize}
\item \textsc{move\_lease}$(f, r)$ makes region $r$ the primary for shard $f$.
It moves no data and is only legal where $r$ already holds a replica.
\item \textsc{move\_replica\_to}$(f, r)$ places a replica of shard $f$ in
region $r$, evicting the existing replica furthest from $r$ so that the
replication factor stays constant. It costs one shard's worth of movement.
\item \textsc{split}$(f)$ divides shard $f$ at its midpoint.
\item \textsc{noop} does nothing.
\end{itemize}

Two properties of this encoding matter. First, every action names a desired end
state rather than an operation, so every action is defined in every state: an
action that cannot be applied resolves to a small penalty rather than an
error. This matters because a Deep Q-Network cannot mask illegal actions, so an
action space with illegal members would need machinery that this one avoids.
Second, the encoding was chosen by measurement rather than by taste. The
destination-only form of \textsc{move\_replica\_to} gives 261 actions against
861 for the naive source-and-destination form, with a higher fraction of
actions applicable in a typical state, and a breadth-first search over the
reachability graph confirms that no placement becomes unreachable.

\paragraph{Reward} The reward is the negative weighted sum of three normalized
terms: mean latency divided by a scale constant, data moved as a fraction of
the total dataset, and load imbalance. Expressing movement as a fraction of the
dataset rather than in bytes is deliberate and necessary: the simulator holds
60\,GB and the real cluster used for validation holds about 1\,GB, and a reward
denominated in bytes would score the same physical move sixty times
differently in the two environments, making a policy pre-trained in one
meaningless in the other.

\subsection{Baselines}

Four reference points. Three are policies and one is a bound.

\begin{itemize}
\item \textbf{Static hash} never moves anything. The initial placement is what
a consistent-hash partitioner~\cite{karger1997consistent} produces, and it is
held for the whole episode.
This is what a production database does today and it is the thing to beat.
\item \textbf{Random} selects a legal action uniformly. It is a control: it
proves that any gain comes from choosing \emph{where} to move data rather than
from moving data at all.
\item \textbf{Follow traffic} finds, each window, the single shard whose
primary placement costs the most relative to the best available region, and
moves it towards that region. It is the obvious rule an operator would write by
hand, and it is myopic by construction, weighing no migration cost and
anticipating no drift.
\item \textbf{Best per window (fixed shards)} is not a policy. It is the lowest
latency reachable by relocating every primary for free and instantly, each
window, at the shard boundaries the episode started with. No real policy can
reach it, because it pays no movement cost and sees the window it is
optimizing. It is also not a global optimum, because it never splits a shard, so
a policy that splits can legitimately sit below it.
\end{itemize}

\subsection{The agent}

The agent is a Deep Q-Network. A Q-network is a function that, given the
current state, estimates the long-run value of each available action; the
policy is then to take the action with the highest estimate. The network here
has two hidden layers of 256 units and is trained with the standard
Stable-Baselines3 implementation~\cite{raffin2021sb3} against a Gymnasium
environment~\cite{towers2025gymnasium}, following Mnih et
al.~\cite{mnih2015dqn}.

The plain form is used deliberately, without the extensions that
Rainbow~\cite{hessel2018rainbow} combines: double Q-learning to correct
overestimation~\cite{vanhasselt2016double}, a dueling
head~\cite{wang2016dueling}, or prioritized replay~\cite{schaul2016per}. The
question here is whether the problem is learnable at all, and establishing that
with the simplest agent makes any later improvement attributable to the
extension rather than to the combination. The evaluation reports where that
choice costs something.

Two choices in the training setup are worth stating because they are easy to
get wrong. Training resets onto a rotating set of 64 workload seeds rather than
one, because an agent trained on a single realization learns that realization
rather than the structure behind it. Evaluation uses a block of seeds disjoint
from the training seeds and acts greedily rather than exploring, so the
reported number is generalization rather than recall.

\subsection{Metrics}

Every result in this paper reports some subset of five quantities.

\begin{itemize}
\item \textbf{Mean latency}, in milliseconds, averaged over the windows of an
episode.
\item \textbf{p95 latency}, the value below which 95\% of observations fall. It
is reported alongside the mean because an average can hide a bad tail, and
users experience the tail.
\item \textbf{Data moved}, in gigabytes per episode. This is the price paid for
whatever latency was bought.
\item \textbf{Cross-region transaction fraction}, the share of transactions
whose shards span more than one region and which therefore pay the two-phase
commit charge.
\item \textbf{Load imbalance}, a measure of how unevenly requests are spread
across regions, included because a policy can lower latency by piling all the
work onto one region, which is not a real improvement.
\end{itemize}

Results are reported as the mean over eight seeds with a 95\% percentile
bootstrap~\cite{efron1979bootstrap} confidence interval computed over those
seeds. The bootstrap is used rather than a $t$ interval because eight is a small
sample and nothing guarantees the per-seed means are normally distributed.

\subsection{Experimental apparatus}

Every number reported here comes from a single script that writes one
self-contained directory per run, holding the resolved configuration, the seed
list, the commit the code was at, the package versions, the machine, the raw
per-window record, the per-seed summary, and the exact command that reproduces
the run. Runs execute on GitHub Actions rather than on a personal machine, so
each is unattended and each records the hardware it received. This is reported
because it is part of the method: a number that cannot be traced to the
configuration that produced it cannot be defended.

\subsection{Validation on a real cluster}

The structure of pre-training in a cheap simulator and fine-tuning on expensive
real hardware is borrowed from robotics, where it is standard and where the gap
between the two is the central difficulty~\cite{tobin2017domain}. The difference
here is that the gap is narrower and more checkable: the simulator and the
cluster can be made to agree on the observation space, the action space and the
cost structure by construction, which is why the database was chosen for that
property.

The simulator's cost model is an assumption, and the assumption is testable.
I deploy six CockroachDB~\cite{taft2020cockroachdb} nodes, one per simulated
region, in containers on a single host, with wide-area round-trip times injected
between them using the Linux kernel's traffic control facility. CockroachDB was
selected over sixteen other candidate systems on one criterion: it exposes a \emph{synchronous}
statement for relocating a shard's lease and its replicas. Synchrony is
essential rather than convenient. If a placement command returns before the
placement has changed, the reward observed after an action belongs to a state
the action did not produce, and the decision problem is no longer well defined.

\begin{figure}[t]
\centering
\begin{tikzpicture}[
  node distance=3mm,
  box/.style={draw, rounded corners=1pt, align=center, font=\scriptsize,
              minimum width=30mm, inner sep=1.6mm},
  title/.style={font=\scriptsize\bfseries},
  parity/.style={<->, >={Stealth[length=1.6mm]}, dashed, gray},
  note/.style={font=\tiny, gray, align=center}
]

% Left column: the simulator.
\node[title] (simtitle) {Simulator};
\node[box, below=2mm of simtitle] (simwork) {drifting workload\\256 cells, 6 regions};
\node[box, below=of simwork] (simcost) {analytic cost model\\reads, writes, 2PC};
\node[box, below=of simcost] (simobs) {observation encoder};

% Right column: the real cluster.
\node[title, right=26mm of simtitle] (clutitle) {CockroachDB cluster};
\node[box, below=2mm of clutitle] (cluwork) {load generator\\6 nodes, netem latency};
\node[box, below=of cluwork] (clucost) {measured latency\\from real transactions};
\node[box, below=of clucost] (cluobs) {observation encoder};

% The three parity constraints.
\draw[parity] (simwork) -- node[note, above] {action space} node[note, below] {261 actions} (cluwork);
\draw[parity] (simcost) -- node[note, above] {cost structure} (clucost);
\draw[parity] (simobs) -- node[note, above] {observation space} (cluobs);

% The agent, underneath, moving left to right.
\node[box, below=6mm of $(simobs)!0.5!(cluobs)$, minimum width=62mm]
  (agent) {Deep Q-Network\\pre-trained on the left, fine-tuned on the right};
\draw[->, >={Stealth[length=1.6mm]}] (simobs) -- (agent);
\draw[->, >={Stealth[length=1.6mm]}] (cluobs) -- (agent);

\end{tikzpicture}
\caption{The two halves of the system and the three things that must agree
between them. The agent is pre-trained against the simulator, which is fast and
cheap, and fine-tuned against the cluster, which is slow and real. Transfer is
only meaningful if the observation space, the action space and the cost
structure are the same on both sides, which is why the database was chosen for
its ability to satisfy them rather than for its popularity.}
\label{fig:architecture}
\end{figure}
